% !TEX root = ../main.tex
\chapter{合成可能な設計を作る}
\label{chap:composable-design}

\begin{chapterabstract}
本章では、合成可能な設計を、圏論の「射の合成」とLean 4の小さな等式証明を使って整理する。対象は、関数パイプライン、middleware、ETL、adapter chain のように、複数の処理を順に接続する設計である。個々の部品が正しいだけでは、全体が安全につながるとは限らない。重要なのは、各部品の入力型と出力型が明確であり、恒等処理を挿入しても意味が変わらず、括弧の付け方を変えても結果が変わらないことである。

本章では、ETLパイプラインを題材に、処理ステージを関数として表す。まず、合成できる部品とできない部品の違いを確認する。次に、入出力型をそろえる設計原則を説明する。最後に、Leanで合成の結合律と恒等律を小さく証明し、設計レビューで何を確認すべきかを整理する。
\end{chapterabstract}

\begin{learninggoals}
\item 合成可能な設計を、部品の良し悪しではなく、接続可能性の観点から説明できる。
\item 入力型と出力型を明示して、処理ステージの境界を設計できる。
\item 恒等律と結合律を、パイプライン設計レビューの確認項目として使える。
\item ETLパイプラインの小さなLeanモデルを書き、合成後の振る舞いを確認できる。
\end{learninggoals}

\section{この章を読む理由}

前章までで、データ移行、API adapter、DBスキーマ変更のように、比較的はっきりした変換を扱ってきた。本章からは、もう少し広い設計の話へ進む。実務では、最初から「この関数とこの関数が同型である」と言える場面ばかりではない。むしろ、設計段階で「後から安全につなげられる形にしておく」ことが重要になる。

たとえば、ETLでは、生データを取り込み、検証し、補完し、集計し、レポート用の形へ変換する。各処理を一つの巨大な関数にまとめることもできる。しかし、それでは途中の仕様を確認しにくい。逆に、処理を小さく分けても、入出力が曖昧なら安全につなげられない。

\begin{intuitionbox}
合成可能な設計とは、「部品をただ分割する設計」ではない。部品をつないだときに、全体の意味が予測できる設計である。圏論の言葉では、部品を射、つなぐことを合成として見る。Leanでは、その合成が期待通りに振る舞うことを小さな等式として確認できる。
\end{intuitionbox}

\FigurePlaceholder{fig-ch33-overview.png}{0.90}{合成可能な設計}{fig:ch33-overview}

図\ref{fig:ch33-overview}は、型を境界として処理を接続し、個々の部品からパイプラインを組み立てる関係を示している。

\section{合成できる部品とできない部品}

ソフトウェア設計で「部品化」と言うと、関数やクラスやサービスを小さく分けることを思い浮かべやすい。しかし、小さいだけでは合成可能とは言えない。合成可能な部品には、少なくとも次の性質が必要である。

\begin{definitionbox}
本章では、処理部品を「ある型の値を受け取り、別の型の値を返すもの」として扱う。型 \texttt{A} から型 \texttt{B} への処理を、Leanでは \texttt{A -> B} と書く。圏論の言葉では、これは \texttt{A} から \texttt{B} への射として読む。
\end{definitionbox}

部品 \texttt{f : A -> B} と \texttt{g : B -> C} があるとき、\texttt{f} の出力型と \texttt{g} の入力型が一致しているので、\texttt{g \textbackslash{}circ f} として合成できる。ソフトウェア的には、\texttt{f} の結果をそのまま \texttt{g} に渡せるということである。

\begin{longtable}{p{0.20\linewidth}p{0.35\linewidth}p{0.35\linewidth}}
\toprule
観点 & 合成しやすい部品 & 合成しにくい部品 \\
\midrule
境界 & 入力型と出力型が明確 & 引数の意味がコメントにしかない \\
副作用 & 副作用がない、または文脈型に閉じ込めている & ログ、DB更新、外部API呼び出しが内部に混在する \\
エラー & \texttt{Option} や \texttt{Except} などで失敗を型に出す & 例外、null、ログだけで失敗を表す \\
責務 & 検証、補完、集計などが分かれている & 一つの関数が検証、変換、保存を同時に行う \\
法則 & 恒等処理や結合の確認がしやすい & つなぐ順序で隠れた状態が変わる \\
\bottomrule
\end{longtable}

\begin{softwarebox}
「合成できない」ことは、単に型エラーになるという意味だけではない。型としてはつながっても、意味としてつながらない場合がある。たとえば、\texttt{String -> String} の関数は何でも合成できるように見える。しかし、一方が日付文字列を期待し、もう一方がメールアドレスを返すなら、型は同じでも設計上は合成できていない。
\end{softwarebox}

したがって、合成可能性を設計レビューで見るときは、型が合うことと、意味が合うことを分けて確認する必要がある。Leanは前者を強く助ける。後者については、型名、構造体、不変条件、定理名を使って、意味の境界をできるだけ明示する。

\section{入出力型をそろえる}

合成可能な設計の第一歩は、処理ステージごとに入力型と出力型を決めることである。ETLを例にすると、すべてを \texttt{String} や \texttt{Json} のまま流すのではなく、段階ごとに型を変える。

\begin{softwarebox}
たとえば、次のように段階を分ける。
\begin{itemize}
\item \texttt{RawRow}: 外部から取り込んだ生データ。
\item \texttt{ValidRow}: 最低限の検証を通ったデータ。
\item \texttt{EnrichedRow}: 補完や計算済み項目を持つデータ。
\item \texttt{ReportRow}: レポート出力に必要な形へ整えたデータ。
\end{itemize}
このように型名で段階を分けると、「どの処理がどの段階のデータを受け取るか」がレビューしやすくなる。
\end{softwarebox}

\FigurePlaceholder{fig-ch33-boundaries.png}{0.90}{処理ステージの境界}{fig:ch33-boundaries}

図\ref{fig:ch33-boundaries}のように、前段の出力型を後段の入力型に合わせると、接続できるステージが型から分かる。ここで重要なのは、型を細かくしすぎることではなく、意味が変わる境界に型名を与えることである。たとえば、単なる整形処理で意味が変わらないなら、型を分ける必要はないかもしれない。一方、生データから検証済みデータへ移る境界は、仕様上の意味が大きく変わるので、型を分ける価値がある。

\begin{warningbox}
型をそろえることと、すべてを同じ型にすることは逆である。すべてを \texttt{String -> String} にすると、表面上は何でも合成できる。しかし、どの段階のデータなのかを型が教えてくれない。合成可能な設計では、むしろ必要なところで型を分ける。
\end{warningbox}

\section{合成法則を設計レビューで使う}

圏論で最初に出てくる法則は、恒等律と結合律である。抽象的に聞こえるが、設計レビューではかなり実用的な観点になる。

\begin{definitionbox}
恒等律とは、何もしない処理を前後に入れても結果が変わらないという法則である。結合律とは、三つの処理を順につなぐとき、\texttt{(h \textbackslash{}circ g) \textbackslash{}circ f} と \texttt{h \textbackslash{}circ (g \textbackslash{}circ f)} の結果が同じになるという法則である。
\end{definitionbox}

ソフトウェア設計での読み替えは次の通りである。

\begin{itemize}
\item 恒等律: ログだけの wrapper や空の変換ステップを入れても、ドメイン値の意味が変わらない。
\item 結合律: パイプラインを前半と後半に分けても、まとめて実行しても、純粋な変換部分の結果は変わらない。
\item 型整合: 前段の出力が、後段の入力として意味的にも型的にも受け取れる。
\end{itemize}

Leanでは、まず純粋な関数についてこれらを確認する。次のコードは、任意の型 \texttt{A}, \texttt{B}, \texttt{C}, \texttt{D} と任意の関数 \texttt{f}, \texttt{g}, \texttt{h} について、合成の基本法則を点ごとの等式として表している。

\begin{leanbox}
ここでは関数そのものの等しさではなく、「任意の入力 \texttt{x} に対して結果が同じ」という形で書く。これは第21章の関数外延性に進む前でも読みやすい形である。
\end{leanbox}

\begin{listing}[htbp]
\begin{minted}{lean4}
namespace Chapter33

variable {A B C D : Type}

theorem left_id_point (f : A -> B) (x : A) :
    (Function.comp id f) x = f x := by
  rfl

theorem right_id_point (f : A -> B) (x : A) :
    (Function.comp f id) x = f x := by
  rfl

theorem assoc_point
    (f : A -> B) (g : B -> C) (h : C -> D) (x : A) :
    (Function.comp h (Function.comp g f)) x =
      (Function.comp (Function.comp h g) f) x := by
  rfl

end Chapter33
\end{minted}
\caption{\texttt{left\_id\_point}・\texttt{right\_id\_point}・\texttt{assoc\_point}}
\label{lst:ch33-composition-laws}
\end{listing}

\texttt{variable \{A B C D : Type\}} は、この名前空間内の定理で共通して使う型変数をまとめて宣言している。この証明が短いのは、Leanにとって関数合成の両辺が定義を展開すると同じ形になるからである。これは実務で「どんな処理でも順序を変えてよい」という意味ではない。ここで証明しているのは、純粋な関数合成の括弧の付け方だけである。副作用や失敗を含む処理については、\texttt{Option}、\texttt{Except}、状態などの文脈を型に出す必要がある。

\section{サンプル：ETLパイプラインの段階的設計}

ここでは、小さなETLパイプラインを作る。現実のETLでは、ファイル入力、パースエラー、DB更新、外部API、時刻、監査ログなどが関わる。本章では、それらをいったん外に置き、純粋なデータ変換の核だけを取り出す。

\FigurePlaceholder{fig-ch33-etl-pipeline.png}{0.92}{ETLパイプライン}{fig:ch33-etl-pipeline}

図\ref{fig:ch33-etl-pipeline}では、取り込み後の行を検証、補完、集計の順に変換する流れを、型の異なる三つの関数として表している。最初に、段階ごとの型を定義する。ここでは説明を短くするため、金額のパースや日付処理は省略し、金額を最初から \texttt{Nat} で扱う。

\begin{listing}[htbp]
\begin{minted}{lean4}
namespace Ch33ETL

structure RawRow where
  userId : Nat
  rawAmount : Nat
  deriving Repr, DecidableEq

structure ValidRow where
  userId : Nat
  amount : Nat
  deriving Repr, DecidableEq

structure EnrichedRow where
  userId : Nat
  amount : Nat
  fee : Nat
  deriving Repr, DecidableEq

structure ReportRow where
  userId : Nat
  total : Nat
  deriving Repr, DecidableEq
\end{minted}
\caption{\texttt{RawRow}・\texttt{ValidRow}・\texttt{EnrichedRow}・\texttt{ReportRow}}
\label{lst:ch33-etl-types}
\end{listing}

この段階では、まだ圏論らしい抽象は出てこない。ただし、型を分けたことで、\texttt{RawRow -> ValidRow}、\texttt{ValidRow -> EnrichedRow}、\texttt{EnrichedRow -> ReportRow} という射の列が見えるようになる。なお、この簡略モデルの \texttt{ValidRow} は検証済みという不変条件をデータとして保持していない。型名だけで検証を証明したことにはならず、現実の検証をモデル化するなら、条件を満たす値だけを構成できる型にするか、\texttt{RawRow -> Option ValidRow} のように失敗を型へ含める必要がある。

\begin{listing}[htbp]
\begin{minted}{lean4}
def sanitize (r : RawRow) : ValidRow :=
  { userId := r.userId, amount := r.rawAmount }

def enrich (r : ValidRow) : EnrichedRow :=
  { userId := r.userId, amount := r.amount, fee := 10 }

def summarize (r : EnrichedRow) : ReportRow :=
  { userId := r.userId, total := r.amount + r.fee }

def pipeline : RawRow -> ReportRow :=
  Function.comp summarize (Function.comp enrich sanitize)

def pipelineExpanded (r : RawRow) : ReportRow :=
  summarize (enrich (sanitize r))
\end{minted}
\caption{\texttt{sanitize}・\texttt{enrich}・\texttt{summarize}}
\label{lst:ch33-etl-functions}
\end{listing}

\texttt{pipeline} は、三つの処理を合成したものである。\texttt{pipelineExpanded} は、同じ処理を入れ子の関数適用として書いたものである。この例の \texttt{sanitize} は \texttt{Nat} のフィールドを移すだけであり、パースや妥当性検査を実行しない。設計レビューでは、このモデル上の処理内容と、名前から連想される本番処理を区別したうえで、二つの書き方の意味が同じかを確認する。

\begin{listing}[htbp]
\begin{minted}{lean4}
theorem pipeline_same_as_expanded (r : RawRow) :
    pipeline r = pipelineExpanded r := by
  rfl

def pipelineWithIdentity (r : RawRow) : ReportRow :=
  summarize ((fun x => x) (enrich (sanitize r)))

theorem identity_step_removed (r : RawRow) :
    pipelineWithIdentity r = pipeline r := by
  rfl

theorem pipeline_assoc (r : RawRow) :
    (Function.comp summarize (Function.comp enrich sanitize)) r =
      (Function.comp (Function.comp summarize enrich) sanitize) r := by
  rfl

end Ch33ETL
\end{minted}
\caption{\texttt{pipeline\_same\_as\_expanded}・\texttt{identity\_step\_removed}・\texttt{pipeline\_assoc}}
\label{lst:ch33-etl-proofs}
\end{listing}

これらの定理は、現実のETL全体を証明しているわけではない。しかし、純粋な変換部分について、次のような設計意図を検査可能にしている。

\begin{itemize}
\item ステージを合成で書いても、展開して書いても同じである。
\item 何もしない中間ステップを入れても、結果は変わらない。
\item 三つのステージをどこで括っても、純粋な関数合成としての結果は変わらない。
\end{itemize}

\section{ステージという小さな設計単位}

実務では、単なる関数 \texttt{A -> B} だけでなく、処理名、監査情報、設定、メトリクスなどを一緒に持たせたいことがある。その場合でも、核となる変換を \texttt{run : A -> B} として取り出しておくと、合成の議論を保ちやすい。

次のコードでは、\texttt{Stage A B} を「\texttt{A} を受け取り \texttt{B} を返す処理ステージ」として定義する。\texttt{then} は、前のステージの出力を次のステージへ渡す合成である。

\begin{listing}[htbp]
\begin{minted}{lean4}
namespace Ch33Stage

structure Stage (A B : Type) where
  run : A -> B

namespace Stage

variable {A B C D : Type}

def id (A : Type) : Stage A A :=
  { run := fun x => x }

def andThen (s1 : Stage A B) (s2 : Stage B C) : Stage A C :=
  { run := Function.comp s2.run s1.run }

theorem left_id (s : Stage A B) (x : A) :
    ((id A).andThen s).run x = s.run x := by
  rfl
\end{minted}
\caption{\texttt{Stage.id}・\texttt{Stage.andThen}・\texttt{Stage.left\_id}}
\label{lst:ch33-stage}
\end{listing}

このコードでは、処理ステージを\texttt{run : A -> B}を持つ構造体として包み、\texttt{andThen}で二つのステージを合成している。
\texttt{left\_id}は、左側に何もしないステージを置いても実行結果が変わらないことを確認している。

\begin{listing}[htbp]
\begin{minted}{lean4}
theorem right_id (s : Stage A B) (x : A) :
    (s.andThen (id B)).run x = s.run x := by
  rfl

theorem assoc
    (s1 : Stage A B) (s2 : Stage B C) (s3 : Stage C D) (x : A) :
    ((s1.andThen s2).andThen s3).run x =
      (s1.andThen (s2.andThen s3)).run x := by
  rfl

end Stage
end Ch33Stage
\end{minted}
\caption{\texttt{Stage.right\_id}・\texttt{Stage.assoc}}
\label{lst:ch33-stage-laws}
\end{listing}

\texttt{left\_id}、\texttt{right\_id}、\texttt{assoc} が示すのは、任意の入力に対する \texttt{run} の結果の等しさであり、定理の主張そのものは \texttt{Stage} 構造体同士の等式ではない。本章で観測したいのは処理結果なので、この点ごとの形で十分である。\texttt{Stage} は本格的な圏論ライブラリではなく、設計上の処理ステージを小さくモデル化したものにすぎないが、この範囲の恒等律と結合律はLeanに検査させられる。

\section{合成できない処理を合成可能に近づける}

既存の設計が最初から合成可能になっているとは限らない。むしろ、業務ロジック、永続化、ログ、例外処理、外部API呼び出しが一つの関数に混ざっていることの方が多い。その場合は、いきなり全体をLeanに移すのではなく、次のように分解する。

\begin{softwarebox}
合成可能に近づける典型的な分解は次の通りである。
\begin{enumerate}
\item 外部入力を受け取る部分と、純粋な変換部分を分ける。
\item 検証済みデータには、\texttt{ValidRow} のような別の型名を与える。
\item 失敗する処理は、戻り値の型に \texttt{Option} や \texttt{Except} を出す。
\item DB保存や外部API呼び出しは、純粋な変換の後ろに置く。
\item 設計レビューでは、まず純粋な核について合成法則を確認する。
\end{enumerate}
\end{softwarebox}

この分解は、アーキテクチャ上の層分けとも相性がよい。ドメイン層では \texttt{A -> B} に近い純粋な変換を増やす。インフラ層では、外部システムとのやり取りを担う。両者の境界には、型と変換関数を置く。これにより、どの部分をLeanで証明し、どの部分を統合テストや監視で確認するかを分けやすくなる。

\begin{warningbox}
副作用を含む処理でも、何らかの形で合成できることは多い。ただし、その場合の法則は単純な関数合成の法則とは限らない。失敗、状態、非同期、ログなどを含む処理は、第17章のモナドや第34章のeffectful codeのリファクタリング規則と合わせて考える必要がある。
\end{warningbox}

\section{設計レビューで見るチェックリスト}

合成可能な設計をレビューするときは、「部品が小さいか」だけを見るのでは不十分である。次の観点で、部品の接続が安全かを確認する。

\begin{softwarebox}
\begin{itemize}
\item 各ステージの入力型と出力型は明示されているか。
\item 型名は、そのデータがどの段階にあるかを表しているか。
\item 何もしないステップを入れても意味が変わらないか。
\item パイプラインを前半と後半に分けても、純粋な変換結果は変わらないか。
\item 失敗や副作用は、型または境界で明示されているか。
\item 中間データの不変条件は、型名や述語として表せるか。
\item Leanで確認する純粋な核と、通常のテストで確認する外部挙動が分かれているか。
\end{itemize}
\end{softwarebox}

このチェックリストは、圏論の言葉をそのまま設計レビューに持ち込むためのものではない。むしろ、合成、恒等、結合という圏論の最小語彙を、日常のレビュー項目へ翻訳したものである。

\section{Leanで証明したことと実務上の限界}

本章のLeanコードで示したことは、純粋な関数としてモデル化した処理に関する性質である。したがって、次のようなことは本章の証明だけでは保証していない。

\begin{warningbox}
本章の証明は、実DBのトランザクション、ネットワーク障害、外部APIの実挙動、時刻依存、性能、権限設定、監査ログの完全性を保証しない。Leanで証明したのは、モデル化した純粋な変換の核についてである。
\end{warningbox}

それでも、この小さな証明には実務上の価値がある。なぜなら、設計上の重要な判断を明確にできるからである。「このステージはどの型からどの型への変換なのか」「途中に何もしない処理を入れてよいのか」「ステージの括り方を変えても意味が変わらないのか」。これらの問いに答えられる設計は、変更に強い。

\section{本章のまとめ}

合成可能な設計では、部品そのものだけでなく、部品の接続可能性を見る。

入出力型を明示すると、処理ステージの境界と意味がレビューしやすくなる。

恒等律は、何もしない処理を挿入しても意味が変わらないこととして読める。

結合律は、パイプラインの括り方を変えても純粋な変換結果が変わらないこととして読める。

Leanでは、ETLのような実務例から純粋な変換の核を取り出し、小さな等式として検査できる。

副作用、失敗、外部システムは、本章の純粋なモデルとは分けて扱う必要がある。
